A possible source of background for this analysis are muons 
  that arise from decay of pions and Kaons, as well as muons
  arising from secondary interactions in the calorimeter.
In this section, an estimate of the contributions of these
  background pairs to the measurement is made. 
The contribution of the background muons can be statistically 
  estimated using the momentum-imbalance as a discriminant.
The momentum-imbalance for a muon is defined as:


\begin{eqnarray}
  \frac{\Delta p}{p_{\mathrm{ID}}}=\frac{(p_{\mathrm{ID}}-p_{\mathrm{MS}})}{p_{\mathrm{ID}}}
\end{eqnarray}


where, $p_{\mathrm{ID}}$ is the momentum of the Inner Detector track associated with the muon,
  and $p_{\mathrm{MS}}$ is the Muon-Spectrometer extrapolated track particle momentum,
  which is obtained from extrapolating the Muon-Spectrometer segment of the muon track
  to the interaction region while correcting for the average energy lost by the muon in the calorimeter.
For muons coming from HF decays this distribution is expected to be Gaussian and symmetric about
  $\Delta p/p_{\mathrm{ID}}=0$, however for muons produced via inflight decays such as pions and kaons, which become
  misidentified as muons, the  $\Delta p/p_{\mathrm{ID}}$ is much broader and shifted to a positive mean value
  of $\Delta p/p_{\mathrm{ID}}$ (i.e. $\langle\Delta p/p_{\mathrm{ID}}\rangle>0$).
Cuts on this quantity can be used to reduce the background particles identified as muons. 
However, both the mean and the width of this quantity can be highly $\eta, \pt,$ 
  and centrality dependent, as shown in Figure~\ref{fig:Dpopmeanandrms}.
Therefore, a simple cut on this quantity would introduce \pt\ and $\eta$ dependent
  efficiencies.
For this reason, it became pertinent to construct a $\Delta p/p$ significance, 
  defined as
\begin{eqnarray}
  \left(\frac{\Delta p}{p}\right)^{\mathrm{sig}}=\frac{(\Delta p/p- \langle\Delta p/p\rangle)}{\sigma(\Delta p/p)}.
  \label{eq:dpopsig}
\end{eqnarray}
where  $\langle\Delta p/p\rangle$ and $\sigma(\Delta p/p)$ are the 
  mean and standard-deviation of the momentum-imbalance distributions
  for true muons obtained from MC, evaluated in fine intervals of \pt\ 
  and $q*\eta$ (Figure~\ref{fig:Dpopmeanandrms}).
Converting the imbalance to a significance before imposing any cuts on it,
  makes the inefficiency imposed by the cuts to be independent of \pt\ and $\eta$.

%The significance was calculated through a combination of values obtained from both
%  HIJING simulations as well as 2015 and 2018 Pb+Pb data.


\begin{figure}
\centering
\includegraphics[width=1.0\textwidth]%                                                                                           
{figures/DpopAna/dpopmeanandrms}
\caption{
Mean (left) and RMS (right) of the momentim-imbalance for (true) muons, 
  as a function of \pt\ for several different $q*\eta$ ranges.
The values are obtained from the HIJING-overlay sample (Section~\ref{sec:PbPbMc}).
\label{fig:Dpopmeanandrms}
}
\end{figure}




Once the single-muon significances are obtained (Eq.~\eqref{eq:dpopsig}),
  a \textit{pair-significance} is obtained for the muon pair as, the
  quadrature-sum of the individual significances:
\begin{eqnarray}
  \left(\frac{\Delta p}{p}\right)^{\mathrm{sig-pair}}=\
  \sqrt{\left(\left(\frac{\Delta p}{p}\right)^{\mathrm{sig-a}}\right)^2 +\
        \left(\left(\frac{\Delta p}{p}\right)^{\mathrm{sig-b}}\right)^2}.
  \label{eq:dpopsigpair}
\end{eqnarray}
To estimate the number of background pairs, the 
  $\left(\frac{\Delta p}{p}\right)^{\mathrm{sig-pair}}$ distributions
  measured in the data are fit using a three-component template-fit,
  with the following components obtained from MC:
\begin{enumerate}
  \item \textbf{Signal-Signal template}: 
    Distribution for $\left(\frac{\Delta p}{p}\right)^{\mathrm{sig-pair}}$ 
      when both muons are signal muons.
  \item \textbf{Signal-Background template}:
    Distribution for $\left(\frac{\Delta p}{p}\right)^{\mathrm{sig-pair}}$ 
      when one muon is a signal muon, and the other is a background muon.
  \item \textbf{Background-Background template}:
    Distribution for $\left(\frac{\Delta p}{p}\right)^{\mathrm{sig-pair}}$ 
      when both muons are background muons.
\end{enumerate}
Replacing the momentum-imbalance distributions by the significances means
  that all of the above distributions become independent of \pt\ and $\eta$
  of the two muons in the pair.
Thus instead of performing the analysis in narrow bins of $(\pta,\ptb,q^a*\etaa,q^b*\etab)$
  we can freely perform the template-analysis across unrestricted regions of
  $(\pta,\ptb,q^a*\etaa,q^b*\etab)$.


Figure~\ref{fig:Templatefits_bycent_qual1_pt4p0to15p0_sign2} shows the template-fit 
  to the pair-significance distributions in the data using the MC-templates.
The plot is for Tight muon selection, \ptab>4~\GeV, and shows the fits for several centrality intervals.
The contributions of the individual templates to the fit are also shown.
From these fits, the fraction of signal-pairs, i.e. the relative contribution of 
  the \textbf{Signal-Signal} are obtained.
Figure~\ref{fig:Signalfractions_qual1}, shows the \textit{signal-fractions} as a function
  of centrality for Tight muons.
It is observed that the pairs in the data have a purity of more than 98\%.
Therefore such, no additional $\Delta p/p$ cuts are imposed in this analysis.



\begin{figure}
\centering
\includegraphics[width=1.0\textwidth]%                                                                                           
{figures/DpopAna/Templatefits_bycent_qual1_pt4p0to15p0_sign2}
\caption{
Template fits to data (in black) for the $\Delta p/p$ pair significance. 
The   \textbf{Signal-Signal} contribution is shown in green, 
  the \textbf{Signal-Background} in magenta and 
  the \textbf{Background-Background} in blue.
The combined fit is shown in red.
The individual panels correspond to different centrality intervals.
The fits are for \ptab>4~\GeV, for Tight muons, and for the combined-charge pairs.
\label{fig:Templatefits_bycent_qual1_pt4p0to15p0_sign2}
}
\end{figure}





\begin{figure}
\centering
\includegraphics[width=1.0\textwidth]%
{figures/DpopAna/Signalfractions_qual1}
\caption{
Signal fractions obtained from the template fits for same-sign muon pairs (open circles), 
  opposite-sign muon pairs (open squares), and their combination (closed circles) 
  as a function of centrality.
Plots are for Tight muons.
\label{fig:Signalfractions_qual1}
}
\end{figure}



\begin{comment}

Although $\eta$-dependence is evident, one of the first tasks in constructing the significance 
  is ascertaining suitable ranges on which to partition the pseudorapidity, 
  or rather the pseudorapidity multiplied by the charge, $\eta q$ to account 
  for the assymmetry of the detector.
After analyzing the detector geometry, borrowing from past accepted values and 
  analyzing $\eta$ windows in 0.2 gaps, final values were decided to be:
$-2.40<\eta q<-1.50$, $-1.50<\eta q<-1.05$, $-1.05<\eta q<-0.50$, $-0.50<\eta q<-0.10$, $-0.10<\eta q<0.10$, $0.10<\eta q<0.50$, $0.50<\eta q<1.05$, $1.05<\eta q<1.50$, $1.50<\eta q<2.40$.
While the ranges of $\eta q$ need to split into positive and negative regions to 
  account for detector asymmetries, the $-0.10<\eta q<0.10$ range can combined for 
  greater statistics because the primary effects in 
  this region are less asymmetrically dependent.


In all but the $-0.1<\eta q<0.1$ bin, data and HIJING overlay values of the $\Delta p/p$ RMS were consistent.
In this range however, due to physical gaps in the detector, unaccounted for fully in HIJING, some discrepancies were noticed.
As this study is to examine dimuon pairs generated via heavy flavor decays, generally speaking, acoplanarity and asymmetry cuts
  on the pairs are employed to remove any photo-nuclear contamination.
However, one can invert this restriction temporarily to study muon pairs which nearly exclusively come from
  $\gamma\gamma \rightarrow \mu^{+}\mu^{-}$ and obtain a very pure sample of dimuons to compare to those of HIJING simulations.
Figure~\ref{fig:Smalletawindowcomparison} shows this comparison (on the left $\gamma\gamma \rightarrow \mu^{+}\mu^{-}$ candidates, and on the right
  HIJING simulation) in the regions of interest.
  
\begin{figure}
\centering
\includegraphics[width=1.0\textwidth]%                                                                                           
{figures/DpopAna/smalletawindowcomparison}
\caption{
The $\Delta p/p$ RMS of $\gamma\gamma \rightarrow \mu^{+}\mu^{-}$ candidates (left) and HIJING samples (right) for small $\eta$ ranges as a function of $\pt$ for an integrated centrality.
\label{fig:Smalletawindowcomparison}
}
\end{figure}  

We can see that while the data is mostly consistent with HIJING for the RMS, HIJING tends to over-predict the
  width in the $-0.10<\eta q<0.10$ range.
As such, the $-0.10<\eta q<0.10$ RMS is ascertained from data and then used to scale the simulation, as will be explained momentarily.

Additionally, while inspecting this region of the detector as it relates to the $\Delta p/p$ RMS, it was noticed that the
   $\gamma\gamma \rightarrow \mu^{+}\mu^{-}$ mean varied noticeably in most all regions of the detector when compared to HIJING,
   as can be seen in Figure~\ref{fig:Gamgam_hijing_compare}. As such, the $\Delta p/p$ mean is also extracted from $\gamma\gamma \rightarrow \mu^{+}\mu^{-}$ data.
 
\begin{figure}
\centering
\includegraphics[width=1.0\textwidth]%                                                                                           
{figures/DpopAna/gamgam_hijing_compare}
\caption{
REMOVE?
\label{fig:Gamgam_hijing_compare}
}
\end{figure} 
  

%And in looking at Figure~\ref{fig:Dpoprmsbyeta}, we see the necessity for separating the $0.00<|\eta q|<0.10$ and $0.10<|\eta q|<0.50$ ranges,
%  but as the $1.50<|\eta q|<2.00$ and $2.00<|\eta q|<2.40$  ranges remaining consistent in both quantities of interest,
%  we combine these bins in our upcoming parametrization procedure for greater statistics.
  

%While these figures show the absolute value of $\eta q$, the final ranges of $\eta q$ were split into positive and
%With these values decided, the mean and the width can be extracted in each and interpolated for the former and
%  parametrized in the latter.

The means showed no significant centrality dependence and hence the integrated $0-100\%$ centrality was used.
The final $\Delta p/p$ as a function of $\pt$ for an integrated centrality separated by $\eta q$ ranges is shown in Figure~\ref{fig:Dpopmeanbyeta}.
As no consistent $\pt$ parametrization is evident amongst all ranges an interpolation was employed in each $\eta q$ shown.  

\begin{figure}
\centering
\includegraphics[width=1.0\textwidth]%                                                                                           
{figures/DpopAna/dpopmeanbyeta}
\caption{
The $\Delta p/p$ mean of $\gamma\gamma \rightarrow \mu^{+}\mu^{-}$ candidates for an integrated centrality for varying $\eta$ ranges.
\label{fig:Dpopmeanbyeta}
}
\end{figure}



The parametrization of the width proved to be somewhat more involved.
Figure~\ref{fig:Dpoprmsbyeta} shows an example width in the 0-10\% centrality as a function of $\pt$ for the given $\eta q$ ranges.
A linear fit was decided on for the appropriate fit.
However, as there is no reason to believe the RMS should increase with $\pt$, values of the slope were restricted to be less than 0.
The slopes of each range were then extracted and plotted as a function of $\eta q$.
The results with a parabolic fit are plotted in the left-hand side of Figure~\ref{fig:Rmsslopefits}.


\begin{figure}
\centering
\includegraphics[width=1.0\textwidth]%                                                                                           
{figures/DpopAna/dpoprmsbyeta}
\caption{
The $\Delta p/p$ RMS with associated linear fit for HIJING samples for an example 10-20$\%$ centrality for varying $\eta$ ranges.
\label{fig:Dpoprmsbyeta}
}
\end{figure}

\begin{figure}
\centering
\includegraphics[width=1.0\textwidth]%                                                                                           
{figures/DpopAna/rmsslopefits}
\caption{
The extracted $\Delta p/p$ RMS slope of the integrated centrality of varying $\eta$ ranges and the associated parabolic fit (left), 
and the $\Delta p/p$ RMS intercept from said slopes (right).
\label{fig:Rmsslopefits}
}
\end{figure}


This fit is the parametrization used to calculate the width for all $\pt$ and  $\eta q$ ranges.
As a certain amount of centrality dependence was observed in the width, this was retained via the intercepts.
The intercepts for each centrality of a function of $\eta q$ can be seen in the right-hand side of Figure~\ref{fig:Rmsslopefits}.
These bin by bin values were used directly for each centrality and $\eta q$ as shown.
However, as mentioned, there was a rather large discrepancy between the data and simulations in the $-0.10<\eta q<0.10$ range as seen here.
As such, these points shown in Figure~\ref{fig:Rmsslopefits} are shown for HIJING but were not used in the final evaluation.
Instead, in this $\eta q$ range, the integrated centrality slope and intercept were found independently by using the
  $\gamma\gamma \rightarrow \mu^{+}\mu^{-}$ candidates.
The slope found from this was used directly, as it is for the other $\eta$ ranges.
For the intercept however, the integrated centrality of the HIJING simulations was scaled to the integrated centrality of the data.
This fraction was then multiplied to each of the intercept centralities found from the the simulation,
  and this final value was used as the intercept for each centrality for the $-0.10<\eta q<0.10$ region.
As such, the $\Delta p/p_{RMS}$ is then calculated as a function of $\pt$, $\eta q$, and centrality defined by:

\begin{eqnarray}
  \Delta p/p_{RMS}=(4.4x10^{-4}\eta^2-2.2x10^{-3})\pt+b_{\eta,c}
\end{eqnarray}

where $b_{\eta,c}$ represents the intercept for a given $\eta$ and centrality range.
As values of the slope could feasibly rise above 0 at extreme $\eta$, these are manually held below 0 during computation.
Additionally, as muons with exceptionally large $\pt$ could unrealistically alter the RMS, muons with $\pt$>15 GeV are held to the same slope value of muons at $\pt$=15 GeV.


In the left-hand plot of Figure~\ref{fig:Dpop_sig_meancompare} we see the $\Delta p/p$ mean plotted as a function of $\pt$ for each $\eta q$ range.
In the right-hand plot of  Figure~\ref{fig:Dpop_sig_meancompare} we see the $\Delta p/p$ significance mean plotted as a function of $\pt$ for each $\eta q$ range.
As desired, the $\eta q$ and centrality dependence in the mean has been removed by creating the $\Delta p/p$ significance.

\begin{figure}
\centering
\includegraphics[width=1.0\textwidth]%                                                                                           
{figures/DpopAna/dpop_sig_meancompare}
\caption{
The $\Delta p/p$ mean (left) and and the mean significance (right) for integrated centrality and varying $\eta$ ranges for HIJING samples.
\label{fig:Dpop_sig_meancompare}
}
\end{figure}


Additionally, in Figure~\ref{fig:Dpop_sig_rmscompare} we compare the width of $\Delta p/p$ (left) to that of its significance (right), as a function of $\pt$ for each $\eta q$ range.
Again, we see the $\eta q$ and centrality dependence minimize in all zones as designed.

\begin{figure}
\centering
\includegraphics[width=1.0\textwidth]%                                                                                           
{figures/DpopAna/dpop_sig_rmscompare}
\caption{
The $\Delta p/p$ RMS (left) and and the RMS significance (right) for integrated centrality and varying $\eta$ ranges for HIJING samples.
(Note: the deviation seen in the $-0.10<\eta<0.10$ is to be expected as the significance values are obtained from $\gamma\gamma \rightarrow \mu^{+}\mu^{-}$ candidates as explained above.)
\label{fig:Dpop_sig_rmscompare}
}
\end{figure}



\end{comment}
